% -----------------------------------------------------------------------------
% Personal customisations: extra packages, structure slides, helper macros.
% Loaded from main.tex with % -----------------------------------------------------------------------------
% Personal customisations: extra packages, structure slides, helper macros.
% Loaded from main.tex with % -----------------------------------------------------------------------------
% Personal customisations: extra packages, structure slides, helper macros.
% Loaded from main.tex with % -----------------------------------------------------------------------------
% Personal customisations: extra packages, structure slides, helper macros.
% Loaded from main.tex with \input{customize.tex}.
%
% Nothing here may change the theme's fonts.  Loading lmodern (or any other
% font package) after the theme silently overrides the Caladea/Carlito pairing
% the theme sets up, which is why it is not loaded.
% -----------------------------------------------------------------------------

\usepackage{amsmath,amssymb}
\usepackage{bm}
\usepackage{booktabs}
\usepackage{listings}

% -----------------------------------------------------------------------------
% Bibliography (APS-like, citations as slide footnotes)
% -----------------------------------------------------------------------------
\usepackage[backend=biber,style=phys,citestyle=numeric,
            giveninits=true,maxbibnames=3,minbibnames=1,
            articletitle=true]{biblatex}
\addbibresource{refs.bib}

% Footnote citations on slides: small, unobtrusive, no separating rule.
\renewcommand*{\bibfont}{\scriptsize}
\setbeamertemplate{bibliography item}{}
\setbeamerfont{footnote}{size=\tiny}
\setbeamercolor{footnote}{fg=unibasGrey}
\setbeamercolor{footnote mark}{fg=unibasGrey}
\renewcommand*{\footnoterule}{}

% -----------------------------------------------------------------------------
% Structure slides
% -----------------------------------------------------------------------------

% Table of contents: one line per section, subsections hidden.
\setbeamertemplate{section in toc}{$\blacktriangleright$~\inserttocsection}
\setbeamertemplate{section in toc shaded}[default][30]
\setbeamertemplate{subsection in toc}{}

% Frames broken over several slides are not numbered "(cont.)".
\setbeamertemplate{frametitle continuation}{}

% A standalone contents slide, with no section highlighted.  Use it once, after
% the title slide; \AtBeginSection takes over from there.
\newcommand{\outlineslide}{%
  \begingroup
  \setbeamertemplate{footline}{}
  \setbeamertemplate{headline}{}
  \themecolor{mint}
  \begin{frame}[plain]{Outline}
    \tableofcontents
  \end{frame}
  \endgroup
}

% Each section opens with the outline, current section highlighted.
\AtBeginSection[]{%
  \begingroup
  \setbeamertemplate{footline}{}
  \setbeamertemplate{headline}{}
  \themecolor{mint}
  \begin{frame}[plain]{Outline}
    \tableofcontents[currentsection]
  \end{frame}
  \endgroup
}

% Each subsection opens with its own divider slide.
\AtBeginSubsection[]{%
  \begingroup
  \setbeamertemplate{footline}{}
  \setbeamertemplate{headline}{}
  \begin{frame}[plain,c]
    \centering
    {\usebeamerfont{framesubtitle}\usebeamercolor[fg]{framesubtitle}%
     \smallcapsline{Section \thesection.\thesubsection}\par}
    \vspace{1.5ex}
    {\usebeamerfont{title}\usebeamercolor[fg]{title}\subsecname\par}
  \end{frame}
  \endgroup
}

% -----------------------------------------------------------------------------
% Code listings
% -----------------------------------------------------------------------------
\definecolor{codegray}{rgb}{0.5,0.5,0.5}
\definecolor{codepurple}{rgb}{0.58,0,0.82}

\lstdefinestyle{unibasel}{
    commentstyle=\color{unibasPetrol},
    keywordstyle=\color{unibasRed},
    numberstyle=\tiny\color{codegray},
    stringstyle=\color{codepurple},
    basicstyle=\ttfamily\scriptsize,
    breakatwhitespace=false,
    breaklines=true,
    captionpos=b,
    keepspaces=true,
    numbers=left,
    numbersep=5pt,
    showspaces=false,
    showstringspaces=false,
    showtabs=false,
    tabsize=4,
    xleftmargin=10pt,
    xrightmargin=10pt,
}
\lstset{style=unibasel}

% -----------------------------------------------------------------------------
% Helper macros
% -----------------------------------------------------------------------------

% Coloured hyperlinks
\newcommand{\hrefcol}[2]{\textcolor{unibasPetrol}{\href{#1}{#2}}}
\newcommand{\hlinkcol}[1]{\hrefcol{#1}{#1}}

% Centred, narrower paragraph (a statement that should not run edge to edge)
\newcommand{\centerstate}[1]{%
  \centering
  \begin{columns}
    \begin{column}{0.8\textwidth}
      #1
    \end{column}
  \end{columns}%
}

% Coloured emphasis.  \c... is Basel red (the highlight colour),
% \b... is petrol (structure).  Use sparingly: on a slide, one emphasis wins.
\newcommand{\ctextbf}[1]{\textbf{\textcolor{unibasRed}{#1}}}
\newcommand{\btextbf}[1]{\textbf{\textcolor{unibasPetrol}{#1}}}
\newcommand{\ctextsl}[1]{\textsl{\textcolor{unibasRed}{#1}}}
\newcommand{\btextsl}[1]{\textsl{\textcolor{unibasPetrol}{#1}}}
\newcommand{\cemph}[1]{\emph{\textcolor{unibasRed}{#1}}}
\newcommand{\bemph}[1]{\emph{\textcolor{unibasPetrol}{#1}}}

% About / acknowledgement page
\newcommand{\aboutpage}[2]{%
  \begingroup
  \setbeamertemplate{footline}{}
  \setbeamertemplate{headline}{}
  \themecolor{mint}
  \begin{frame}[plain,c]{#1}
    \centering
    \begin{minipage}{0.85\textwidth}
      \usebeamercolor[fg]{normal text}
      \centering #2
    \end{minipage}
  \end{frame}
  \endgroup
}

% Bibliography slide
\newcommand{\bibliographypage}{%
  \begingroup
  \themecolor{mint}
  \begin{frame}[allowframebreaks,plain]{References}
    \printbibliography[heading=none]
  \end{frame}
  \endgroup
}
.
%
% Nothing here may change the theme's fonts.  Loading lmodern (or any other
% font package) after the theme silently overrides the Caladea/Carlito pairing
% the theme sets up, which is why it is not loaded.
% -----------------------------------------------------------------------------

\usepackage{amsmath,amssymb}
\usepackage{bm}
\usepackage{booktabs}
\usepackage{listings}

% -----------------------------------------------------------------------------
% Bibliography (APS-like, citations as slide footnotes)
% -----------------------------------------------------------------------------
\usepackage[backend=biber,style=phys,citestyle=numeric,
            giveninits=true,maxbibnames=3,minbibnames=1,
            articletitle=true]{biblatex}
\addbibresource{refs.bib}

% Footnote citations on slides: small, unobtrusive, no separating rule.
\renewcommand*{\bibfont}{\scriptsize}
\setbeamertemplate{bibliography item}{}
\setbeamerfont{footnote}{size=\tiny}
\setbeamercolor{footnote}{fg=unibasGrey}
\setbeamercolor{footnote mark}{fg=unibasGrey}
\renewcommand*{\footnoterule}{}

% -----------------------------------------------------------------------------
% Structure slides
% -----------------------------------------------------------------------------

% Table of contents: one line per section, subsections hidden.
\setbeamertemplate{section in toc}{$\blacktriangleright$~\inserttocsection}
\setbeamertemplate{section in toc shaded}[default][30]
\setbeamertemplate{subsection in toc}{}

% Frames broken over several slides are not numbered "(cont.)".
\setbeamertemplate{frametitle continuation}{}

% A standalone contents slide, with no section highlighted.  Use it once, after
% the title slide; \AtBeginSection takes over from there.
\newcommand{\outlineslide}{%
  \begingroup
  \setbeamertemplate{footline}{}
  \setbeamertemplate{headline}{}
  \themecolor{mint}
  \begin{frame}[plain]{Outline}
    \tableofcontents
  \end{frame}
  \endgroup
}

% Each section opens with the outline, current section highlighted.
\AtBeginSection[]{%
  \begingroup
  \setbeamertemplate{footline}{}
  \setbeamertemplate{headline}{}
  \themecolor{mint}
  \begin{frame}[plain]{Outline}
    \tableofcontents[currentsection]
  \end{frame}
  \endgroup
}

% Each subsection opens with its own divider slide.
\AtBeginSubsection[]{%
  \begingroup
  \setbeamertemplate{footline}{}
  \setbeamertemplate{headline}{}
  \begin{frame}[plain,c]
    \centering
    {\usebeamerfont{framesubtitle}\usebeamercolor[fg]{framesubtitle}%
     \smallcapsline{Section \thesection.\thesubsection}\par}
    \vspace{1.5ex}
    {\usebeamerfont{title}\usebeamercolor[fg]{title}\subsecname\par}
  \end{frame}
  \endgroup
}

% -----------------------------------------------------------------------------
% Code listings
% -----------------------------------------------------------------------------
\definecolor{codegray}{rgb}{0.5,0.5,0.5}
\definecolor{codepurple}{rgb}{0.58,0,0.82}

\lstdefinestyle{unibasel}{
    commentstyle=\color{unibasPetrol},
    keywordstyle=\color{unibasRed},
    numberstyle=\tiny\color{codegray},
    stringstyle=\color{codepurple},
    basicstyle=\ttfamily\scriptsize,
    breakatwhitespace=false,
    breaklines=true,
    captionpos=b,
    keepspaces=true,
    numbers=left,
    numbersep=5pt,
    showspaces=false,
    showstringspaces=false,
    showtabs=false,
    tabsize=4,
    xleftmargin=10pt,
    xrightmargin=10pt,
}
\lstset{style=unibasel}

% -----------------------------------------------------------------------------
% Helper macros
% -----------------------------------------------------------------------------

% Coloured hyperlinks
\newcommand{\hrefcol}[2]{\textcolor{unibasPetrol}{\href{#1}{#2}}}
\newcommand{\hlinkcol}[1]{\hrefcol{#1}{#1}}

% Centred, narrower paragraph (a statement that should not run edge to edge)
\newcommand{\centerstate}[1]{%
  \centering
  \begin{columns}
    \begin{column}{0.8\textwidth}
      #1
    \end{column}
  \end{columns}%
}

% Coloured emphasis.  \c... is Basel red (the highlight colour),
% \b... is petrol (structure).  Use sparingly: on a slide, one emphasis wins.
\newcommand{\ctextbf}[1]{\textbf{\textcolor{unibasRed}{#1}}}
\newcommand{\btextbf}[1]{\textbf{\textcolor{unibasPetrol}{#1}}}
\newcommand{\ctextsl}[1]{\textsl{\textcolor{unibasRed}{#1}}}
\newcommand{\btextsl}[1]{\textsl{\textcolor{unibasPetrol}{#1}}}
\newcommand{\cemph}[1]{\emph{\textcolor{unibasRed}{#1}}}
\newcommand{\bemph}[1]{\emph{\textcolor{unibasPetrol}{#1}}}

% About / acknowledgement page
\newcommand{\aboutpage}[2]{%
  \begingroup
  \setbeamertemplate{footline}{}
  \setbeamertemplate{headline}{}
  \themecolor{mint}
  \begin{frame}[plain,c]{#1}
    \centering
    \begin{minipage}{0.85\textwidth}
      \usebeamercolor[fg]{normal text}
      \centering #2
    \end{minipage}
  \end{frame}
  \endgroup
}

% Bibliography slide
\newcommand{\bibliographypage}{%
  \begingroup
  \themecolor{mint}
  \begin{frame}[allowframebreaks,plain]{References}
    \printbibliography[heading=none]
  \end{frame}
  \endgroup
}
.
%
% Nothing here may change the theme's fonts.  Loading lmodern (or any other
% font package) after the theme silently overrides the Caladea/Carlito pairing
% the theme sets up, which is why it is not loaded.
% -----------------------------------------------------------------------------

\usepackage{amsmath,amssymb}
\usepackage{bm}
\usepackage{booktabs}
\usepackage{listings}

% -----------------------------------------------------------------------------
% Bibliography (APS-like, citations as slide footnotes)
% -----------------------------------------------------------------------------
\usepackage[backend=biber,style=phys,citestyle=numeric,
            giveninits=true,maxbibnames=3,minbibnames=1,
            articletitle=true]{biblatex}
\addbibresource{refs.bib}

% Footnote citations on slides: small, unobtrusive, no separating rule.
\renewcommand*{\bibfont}{\scriptsize}
\setbeamertemplate{bibliography item}{}
\setbeamerfont{footnote}{size=\tiny}
\setbeamercolor{footnote}{fg=unibasGrey}
\setbeamercolor{footnote mark}{fg=unibasGrey}
\renewcommand*{\footnoterule}{}

% -----------------------------------------------------------------------------
% Structure slides
% -----------------------------------------------------------------------------

% Table of contents: one line per section, subsections hidden.
\setbeamertemplate{section in toc}{$\blacktriangleright$~\inserttocsection}
\setbeamertemplate{section in toc shaded}[default][30]
\setbeamertemplate{subsection in toc}{}

% Frames broken over several slides are not numbered "(cont.)".
\setbeamertemplate{frametitle continuation}{}

% A standalone contents slide, with no section highlighted.  Use it once, after
% the title slide; \AtBeginSection takes over from there.
\newcommand{\outlineslide}{%
  \begingroup
  \setbeamertemplate{footline}{}
  \setbeamertemplate{headline}{}
  \themecolor{mint}
  \begin{frame}[plain]{Outline}
    \tableofcontents
  \end{frame}
  \endgroup
}

% Each section opens with the outline, current section highlighted.
\AtBeginSection[]{%
  \begingroup
  \setbeamertemplate{footline}{}
  \setbeamertemplate{headline}{}
  \themecolor{mint}
  \begin{frame}[plain]{Outline}
    \tableofcontents[currentsection]
  \end{frame}
  \endgroup
}

% Each subsection opens with its own divider slide.
\AtBeginSubsection[]{%
  \begingroup
  \setbeamertemplate{footline}{}
  \setbeamertemplate{headline}{}
  \begin{frame}[plain,c]
    \centering
    {\usebeamerfont{framesubtitle}\usebeamercolor[fg]{framesubtitle}%
     \smallcapsline{Section \thesection.\thesubsection}\par}
    \vspace{1.5ex}
    {\usebeamerfont{title}\usebeamercolor[fg]{title}\subsecname\par}
  \end{frame}
  \endgroup
}

% -----------------------------------------------------------------------------
% Code listings
% -----------------------------------------------------------------------------
\definecolor{codegray}{rgb}{0.5,0.5,0.5}
\definecolor{codepurple}{rgb}{0.58,0,0.82}

\lstdefinestyle{unibasel}{
    commentstyle=\color{unibasPetrol},
    keywordstyle=\color{unibasRed},
    numberstyle=\tiny\color{codegray},
    stringstyle=\color{codepurple},
    basicstyle=\ttfamily\scriptsize,
    breakatwhitespace=false,
    breaklines=true,
    captionpos=b,
    keepspaces=true,
    numbers=left,
    numbersep=5pt,
    showspaces=false,
    showstringspaces=false,
    showtabs=false,
    tabsize=4,
    xleftmargin=10pt,
    xrightmargin=10pt,
}
\lstset{style=unibasel}

% -----------------------------------------------------------------------------
% Helper macros
% -----------------------------------------------------------------------------

% Coloured hyperlinks
\newcommand{\hrefcol}[2]{\textcolor{unibasPetrol}{\href{#1}{#2}}}
\newcommand{\hlinkcol}[1]{\hrefcol{#1}{#1}}

% Centred, narrower paragraph (a statement that should not run edge to edge)
\newcommand{\centerstate}[1]{%
  \centering
  \begin{columns}
    \begin{column}{0.8\textwidth}
      #1
    \end{column}
  \end{columns}%
}

% Coloured emphasis.  \c... is Basel red (the highlight colour),
% \b... is petrol (structure).  Use sparingly: on a slide, one emphasis wins.
\newcommand{\ctextbf}[1]{\textbf{\textcolor{unibasRed}{#1}}}
\newcommand{\btextbf}[1]{\textbf{\textcolor{unibasPetrol}{#1}}}
\newcommand{\ctextsl}[1]{\textsl{\textcolor{unibasRed}{#1}}}
\newcommand{\btextsl}[1]{\textsl{\textcolor{unibasPetrol}{#1}}}
\newcommand{\cemph}[1]{\emph{\textcolor{unibasRed}{#1}}}
\newcommand{\bemph}[1]{\emph{\textcolor{unibasPetrol}{#1}}}

% About / acknowledgement page
\newcommand{\aboutpage}[2]{%
  \begingroup
  \setbeamertemplate{footline}{}
  \setbeamertemplate{headline}{}
  \themecolor{mint}
  \begin{frame}[plain,c]{#1}
    \centering
    \begin{minipage}{0.85\textwidth}
      \usebeamercolor[fg]{normal text}
      \centering #2
    \end{minipage}
  \end{frame}
  \endgroup
}

% Bibliography slide
\newcommand{\bibliographypage}{%
  \begingroup
  \themecolor{mint}
  \begin{frame}[allowframebreaks,plain]{References}
    \printbibliography[heading=none]
  \end{frame}
  \endgroup
}
.
%
% Nothing here may change the theme's fonts.  Loading lmodern (or any other
% font package) after the theme silently overrides the Caladea/Carlito pairing
% the theme sets up, which is why it is not loaded.
% -----------------------------------------------------------------------------

\usepackage{amsmath,amssymb}
\usepackage{bm}
\usepackage{booktabs}
\usepackage{listings}

% -----------------------------------------------------------------------------
% Bibliography (APS-like, citations as slide footnotes)
% -----------------------------------------------------------------------------
\usepackage[backend=biber,style=phys,citestyle=numeric,
            giveninits=true,maxbibnames=3,minbibnames=1,
            articletitle=true]{biblatex}
\addbibresource{refs.bib}

% Footnote citations on slides: small, unobtrusive, no separating rule.
\renewcommand*{\bibfont}{\scriptsize}
\setbeamertemplate{bibliography item}{}
\setbeamerfont{footnote}{size=\tiny}
\setbeamercolor{footnote}{fg=unibasGrey}
\setbeamercolor{footnote mark}{fg=unibasGrey}
\renewcommand*{\footnoterule}{}

% -----------------------------------------------------------------------------
% Structure slides
% -----------------------------------------------------------------------------

% Table of contents: one line per section, subsections hidden.
\setbeamertemplate{section in toc}{$\blacktriangleright$~\inserttocsection}
\setbeamertemplate{section in toc shaded}[default][30]
\setbeamertemplate{subsection in toc}{}

% Frames broken over several slides are not numbered "(cont.)".
\setbeamertemplate{frametitle continuation}{}

% A standalone contents slide, with no section highlighted.  Use it once, after
% the title slide; \AtBeginSection takes over from there.
\newcommand{\outlineslide}{%
  \begingroup
  \setbeamertemplate{footline}{}
  \setbeamertemplate{headline}{}
  \themecolor{mint}
  \begin{frame}[plain]{Outline}
    \tableofcontents
  \end{frame}
  \endgroup
}

% Each section opens with the outline, current section highlighted.
\AtBeginSection[]{%
  \begingroup
  \setbeamertemplate{footline}{}
  \setbeamertemplate{headline}{}
  \themecolor{mint}
  \begin{frame}[plain]{Outline}
    \tableofcontents[currentsection]
  \end{frame}
  \endgroup
}

% Each subsection opens with its own divider slide.
\AtBeginSubsection[]{%
  \begingroup
  \setbeamertemplate{footline}{}
  \setbeamertemplate{headline}{}
  \begin{frame}[plain,c]
    \centering
    {\usebeamerfont{framesubtitle}\usebeamercolor[fg]{framesubtitle}%
     \smallcapsline{Section \thesection.\thesubsection}\par}
    \vspace{1.5ex}
    {\usebeamerfont{title}\usebeamercolor[fg]{title}\subsecname\par}
  \end{frame}
  \endgroup
}

% -----------------------------------------------------------------------------
% Code listings
% -----------------------------------------------------------------------------
\definecolor{codegray}{rgb}{0.5,0.5,0.5}
\definecolor{codepurple}{rgb}{0.58,0,0.82}

\lstdefinestyle{unibasel}{
    commentstyle=\color{unibasPetrol},
    keywordstyle=\color{unibasRed},
    numberstyle=\tiny\color{codegray},
    stringstyle=\color{codepurple},
    basicstyle=\ttfamily\scriptsize,
    breakatwhitespace=false,
    breaklines=true,
    captionpos=b,
    keepspaces=true,
    numbers=left,
    numbersep=5pt,
    showspaces=false,
    showstringspaces=false,
    showtabs=false,
    tabsize=4,
    xleftmargin=10pt,
    xrightmargin=10pt,
}
\lstset{style=unibasel}

% -----------------------------------------------------------------------------
% Helper macros
% -----------------------------------------------------------------------------

% Coloured hyperlinks
\newcommand{\hrefcol}[2]{\textcolor{unibasPetrol}{\href{#1}{#2}}}
\newcommand{\hlinkcol}[1]{\hrefcol{#1}{#1}}

% Centred, narrower paragraph (a statement that should not run edge to edge)
\newcommand{\centerstate}[1]{%
  \centering
  \begin{columns}
    \begin{column}{0.8\textwidth}
      #1
    \end{column}
  \end{columns}%
}

% Coloured emphasis.  \c... is Basel red (the highlight colour),
% \b... is petrol (structure).  Use sparingly: on a slide, one emphasis wins.
\newcommand{\ctextbf}[1]{\textbf{\textcolor{unibasRed}{#1}}}
\newcommand{\btextbf}[1]{\textbf{\textcolor{unibasPetrol}{#1}}}
\newcommand{\ctextsl}[1]{\textsl{\textcolor{unibasRed}{#1}}}
\newcommand{\btextsl}[1]{\textsl{\textcolor{unibasPetrol}{#1}}}
\newcommand{\cemph}[1]{\emph{\textcolor{unibasRed}{#1}}}
\newcommand{\bemph}[1]{\emph{\textcolor{unibasPetrol}{#1}}}

% About / acknowledgement page
\newcommand{\aboutpage}[2]{%
  \begingroup
  \setbeamertemplate{footline}{}
  \setbeamertemplate{headline}{}
  \themecolor{mint}
  \begin{frame}[plain,c]{#1}
    \centering
    \begin{minipage}{0.85\textwidth}
      \usebeamercolor[fg]{normal text}
      \centering #2
    \end{minipage}
  \end{frame}
  \endgroup
}

% Bibliography slide
\newcommand{\bibliographypage}{%
  \begingroup
  \themecolor{mint}
  \begin{frame}[allowframebreaks,plain]{References}
    \printbibliography[heading=none]
  \end{frame}
  \endgroup
}
